\chapter{Not the Weak Link: 
Chris's Place in the Everett Chain}

\noindent\begin{minipage}[t]{\textwidth}
\lettrine{C}{hris} is a bare twig bifurcately tied to Hugh Everett III, the architect of the Many-Worlds Interpretation of quantum mechanics.
\begin{figure}[H]
    \centering
    \includegraphics[width=0.8\linewidth]{Illustrations/vign-sperm.png}
    \caption{That ejaculate from which Chris was explicably spawned}
\end{figure}

\noindent
The spawn of the union of the 956,508th sperm of his father’s 1,927th ejaculate with the 326,057th egg of his mother's brood. A convergence, a hum-drum moment of marked cosmic precision and serendipity.
\end{minipage}


\lettrine{C}{hris} Everett’s becoming was borne of an unbecoming moment improbable as any single person's coming-to-be can be. Each digit in the complex code of DNA whispering of branching universes where different beings emerge from slightly altered initial conditions. In our universe, a specific union of genetic material resulted in this Chris Everett; in others, as his uncle would have us believe, the same building blocks assembled into entirely different forms who would never know of the blessed Everett theorem.

\begin{figure}[H]
\includegraphics[width=1.0\linewidth]{Illustrations/vign-scion.png}
\caption{Chris in younger days doing some scioning of his own}
\end{figure}

As a \textit{scion} of Hugh Everett III, Chris’s lived reality was always a marvel to him. While his academic path diverged from quantum mechanics, the philosophical shadow of his uncle's work looms large. In a family where even the most improbable outcomes seem to manifest, Chris Everett finds himself grappling with his place in the grand tapestry of existence. 

As it happens, today of all days, Chris Everett famous for his  Multivariate Universal Principal Component Analysis (MUPCA) in financial circles and setting up an investment fund that claimed to be front running the front runners had been contemplating an alternative to splines for making predictions from financial data when Emmi intercepted his timeline. That after his girlfriend, Nikki had snapped once too often
%%%%%%%%%%%%%%%%%%
\section*{The Crunch between matters}

Chris broke off another square of the chocolate bar, the crunch of hazelnuts breaking against the shore of silence that wrestled between them.  Nikki, curled into the sofa beside him, barely looked up from her book.  \textit{"Do you have to eat it that loudly?"}  

Chris blinked.  
\textit{"What?"}  

Nikki exhaled sharply, turning a page with what he suspected was unnecessary force.  
\textit{"You're crunching. It's irritating."}  

He stared at her for a moment, half a piece of chocolate still melting on his tongue.  
\textit{"It’s a chocolate bar, Nikki. It has a warning oon the packaging that it contains nuts."}  

\textit{"And yet,"} she said, with clipped precision, \textit{"you could still try not to crunch on it like you're broadcasting to the whole room."}  

Chris chewed slower, dramatically exaggerating his next bite.  

She slammed her book shut.  
\textit{"You know what? Forget it."}  

And that was that.  

She set the book on the coffee table and disappeared into the bedroom. A silent night loomed ahead. Chris sat there, staring at the unfinished chocolate bar in his hand. It was absurd. Mindbogglingly absurd. And yet, this was not the first time.  Some minor infraction. A sigh in the wrong place. A hesitation before responding. A misjudged joke. And then this. The silent treatment, the indignation, the exile over a trifling matter that would barely merit a raised eyebrow in any reasonable world.  

\smallskip
And yet, in this world, it  mattered. Ever the Everettian, Chris then thought of an alternative universe: one where he had grabbed from the shelf the plain chocolate bar instead of the one laced with hazelnuts. Where another version of Chris, crunchless, sat here with Nikki in perfect harmony.  

\smallskip

Somewhere out there, that Chris existed. That Chris never was not a the mercy of a hazelnut breaching the delicate soundscape of an evening. That Chris had a peaceful night ahead of him.  

\smallskip
And yet, was that one any happier?  Chris frowned, setting the chocolate aside.  Yes it was nuts. But it wasn’t the nuts. It was never really the nuts. There was always something, some small friction of personality that inveigled itself into their recent conflicts.  

\smallskip

This hypersensitivity to minor annoyances, was that warning sign. A fundamental misalignment; another in a series of pointless little battles that, in the grand scheme of things, would come to the the nothing that is?  And then, with quiet resignation, he thought:  

\textit{Wouldn't it be nice, just once, to live in the universe where consequences didn't always seemlessly reckon from our choices?}

\begin{figure}[H]
\centering
\includegraphics[width=0.9\linewidth]{Illustrations/vign-chrisWife.png}
\caption{Chris making contacting with home}
\end{figure}

He resolved to grab back the control in his life. And promptly followed Nikki into the bedroom and gave her a silent cuddle.

\newpage

\section*{Jammy Dodgers and Markov Chains}

\lettrine{E}{mmi} had just finished her first week in a junior fund risk management program at BlackRock when she found herself at an LSE\footnote{London School of Economics} seminar on Value at Risk modeling in catastrophic settings. The room was filled with academics, some industry professional types, and a handful of postdocs who were more interested in the free tea and biscuits than the lecture. Among them was Chris, a postdoc in financial mathematics under Tiziana's guidance. This was a catastrophic setting.

Emmi spotted Chris across the room wearing a lonely looking tie possibly marinating on some theory, more likely hoping for someone to talk to him.  She decided to oblige him given that the latter was the more likely and he would be more accommodating than his protective airs conveyed. Emmi, armed with a jammy dodger and her sister’s voice in her ear: "Don’t let the academics intimidate you; charm them", approached him with a polite smile.

\begin{figure}[H]
    \centering
    \includegraphics[width=\linewidth]{Illustrations/vign-emmiSeminar.png}
    \caption{Emmi mentally tossing a coin as to hone in on Chris all tied up but alone.}
    \label{fig:meanRevertDist}
\end{figure}


"You’re Chris, right?", she asked. Not waiting, " Hi, I am Emmi my boss sent me down hear to see what was keeping you academics occupied at night- that is as it were- in a world promising many a road to catastrophe. 

Chris raised an eyebrow, conscious of holding his poise, "You'll be the Blackrock delegation then I guess", scanning her with a mix of disdainful curiosity and skepticism. "BlackRock pays for the best and the brightest and yet they still want to hear from the cardigan classes?"

\subsection*{Mean-Reverting Walks and Transitioning}

Emmi tutting, brushed off the jibe. "True, but understanding the models we use to manage risk isn’t optional. You academics get to theorize in peace; we have clients breathing down our necks."

Chris smirked, impressed by her forthright manner. "Fair enough. Let’s talk Markov chains, then."

Emmi batted that diffidence off without a passing thought, "That bit in the seminar about those mean-reverting processes built from transition matrices within a Markov chain setting an—can you walk me through it? I’m still connecting the dots between theory and practice."

Chris put his cold cup of tea down on the table, instantly regretted that decision, and folded his arms. "A Markov chain is all about transitions—how a system moves between states. The beauty of it lies in its simplicity: the future state depends only on the current state, not the path that got you there. That’s what we call the Markov property." 

He reached for his phone and pulled up some code he had prepared for his class next week and then grabbed a napkin and sketched a quick diagram. "In finance, you can model a mean-reverting process with a Markov chain. Imagine a coin toss where the probabilities of heads and tails are adjusted dynamically to pull the outcome back toward zero. This mean-reversion behavior can be controlled by a parameter \(m\), which determines how strongly the system resists deviation." He jotted down the probabilities strangely written in python it seemed:

{
\tiny % Text size adjustment starts here
\begin{verbatim}
mean_reversion_probabilities(h, n, m):
    p_H = 0.5 * (1 - h / n**(1/m))
    q_T = 1 - p_H
    return p_H, q_T
\end{verbatim}
}

Emmi, nibbling on the poor fare that was her biscuit, leaned in. "So, \(m\) here is like a dial. The higher the \(m\), the more tightly the walk sticks to zero?"

Chris nodded. "Exactly. And by tweaking \(m\), you can simulate different market behaviors. Here’s an example."

\subsection*{Markov Chains and Risk Management}

Emmi pointed to the graph. "I can just about see how this would apply to modeling asset prices or volatility. But where does the \textit{stationarity} of the distribution come into play?"

\begin{figure}[H]
    \centering
    \includegraphics[width=1.0\linewidth]{Illustrations/meanreversionEquilibroumWalks.png}
    \caption{Mean-reverting coin toss probability distributions under varying \(m\).}
    \label{fig:meanRevertDist}
\end{figure}

Chris, touched by the relative openness with which she was revealing some ignorance (compared to his students), explained, "The stationary distribution is where the system stabilizes in the long run. Think of it as the equilibrium point of a Markov chain. For mean-reverting processes, the stationary distribution tells us the likelihood of being in any given state after enough time has passed."

% \subsection*{Winning Each Other Over}

Emmi smiled, her earlier nerves displaced by genuine interest. She was now actually listening: "That’s really fascinating. I guess I’ve been so focused on applications that I’ve overlooked the elegance of the theory."

\smallskip

Chris softened, his initial disdain fading. "And I’ve probably dismissed the practical side too quickly. Models mean nothing without real-world constraints to test them. It’s easy to get lost when we play with our mathematical toys."

\smallskip

They simultaneously guffawed awkwardly, the academic come corporate tension morphing in mutual formalism. Emmi reached for another biscuit, glancing at Chris. "You know, my sister Emily always says I should spend more time learning from people who really live down among the equations. She might be right." Chris, now entertaining such barbs as flirtatious: "And maybe I should spend more time talking to people who are grounded in reality. Maybe you can show me how your crew will be using all this modelling one day before the next catastrophe strikes." A pause in maneuvers then passed between them.

%%%%%%%%%%%%%%%%%

\section*{An ex poste Bayesian Chat over Prime Gaps}

\lettrine{“S}o, Chris,” Emmi punctuated the pause in faux-cordialities, tapping her teacup thoughtfully, “about the blended likelihood approach the speaker mentioned. I get that we’re combining empirical data and theoretical priors, but why did it feel like he was just fitting parameters to get a decent hit rate?”

Chris backing off ever so slightly, brushing crumbs or something off his jacket. “It’s not just fitting parameters, Emmi. It’s about balancing two competing goals: being specific enough to capture local behaviors and general enough to work across different scales. That’s why he blended empirical likelihoods with theoretical models.”

“Ah, yes,” Emmi teased, “because nothing says 'reliable model' like \(\alpha\), the magic blending factor.” She smirked. “How do you even decide its value?”

Chris chuckled. “For any blended model, \(\alpha\) isn’t arbitrary—it reflects how much trust you place in your training data versus theory. In finance, you’d call it weighting between historical volatility and implied volatility. 

Emmi, with Emily rather than her boss in her ear chimed, "If say you were playing with just numbers for numbers sake, and were considering training a model on the gaps between primes given some knowledge that they come along according to Poisson how would you begin to parameterize it?.”

Chris caught her gaze (for too long) then fixed on the plot he had just pulled up to show her. He tapped the table gently, his post-seminar tea long forgotten. 

\begin{figure}[H]
    \centering
    \includegraphics[width=\linewidth]{Illustrations/blended-gap-index5000-100.png}
    \caption{Difference Between Bayesian Blended Predicted and Actual Prime Gaps for a Training Window of 5000 Primes}
    \label{fig:blended_gap_diff}
\end{figure}

“So weird that you should say that Emmi,” he began, pointing to the deviations in the chart, “this is what I was going to show my class next week - a totally spurious but hopefully pedagogically insightful way to view the time series modelling process by looking at first difference frequencies in the primes."

"Here I trained the model on the first 5000 primes, using a scaling factor to dampen the impact once we move outside the training window, a smoothing factor at 0.9 for momentum, and blending factor \(\alpha = 0.1\).”

Emmi half looking, then fully examining the plot closely. “It’s… interesting,” she said hesitantly. “The model does seem to capture something, but the deviations are all over the place. How do you even know those three parameters are meaningful?”

Chris chuckled. “That’s exactly the point, Emmi. There’s too much freedom here-this is the embodiment of model risk. I can tweak the scaling factor, adjust the smoothing, or fiddle with \(\alpha\), and I’ll get a slightly different result each time. But does it tell us anything fundamental about the primes? That’s the real question.”

“So, it’s like backtesting a trading strategy,” Emmi said, catching on. “You use past data to predict future behavior, then evaluate the strategy’s performance on unseen data.”

Chris, impressed. “Exactly. And for primes, the gaps are analogous to log price differences in finance. Both are inherently unpredictable at small scales but exhibit statistical patterns over larger scales.”

Emmi frowning rather over-dramatically, but knowingly tellingly “Okay, but looking closer at its predictions so far—it's static hit rates, with limited improvement—don’t inspire confidence. Shouldn’t the model adapt as more gaps are observed? Like, update itself dynamically?”

Chris nodded. “Exactly. That’s where posterior updates come in. Each observed gap modifies the prior, making the model smarter as it learns. Think of it like Bayesian reasoning—predict, observe, adjust, repeat.”

Chris held grace with his laptop. "To make this concrete, let me show you the core of the code that blends empirical and theoretical likelihoods. This snippet computes the blended likelihood, which is the backbone of the Bayesian model."

% \section*{Blended Likelihood Code}

{
\tiny % Text size adjustment starts here
\begin{verbatim}
def compute_blended_likelihood(train_gaps, lambda_param, alpha):
    # Empirical likelihood from training data
    unique_gaps, gap_counts = np.unique(train_gaps, return_counts=True)
    empirical_likelihood = {gap: count / len(train_gaps) 
    for gap, count in zip(unique_gaps, gap_counts)}
    # Theoretical Poisson likelihood
    max_gap = max(train_gaps)
    poisson_likelihood = {gap: (lambda_param ** gap) * np.exp(-lambda_param) / np.math.factorial(gap) 
    for gap in range(1, max_gap + 1)}
    # Blending the two likelihoods
    blended_likelihood = {}
    for gap in range(1, max_gap + 1):
        emp_likelihood = empirical_likelihood.get(gap, 0)
        poi_likelihood = poisson_likelihood.get(gap, 0)
        blended_likelihood[gap] = alpha * emp_likelihood + (1 - alpha) * poi_likelihood
    # Normalizing the blended likelihood
    total = sum(blended_likelihood.values())
    for gap in blended_likelihood:
        blended_likelihood[gap] /= total
    return blended_likelihood
\end{verbatim}
} % Text size adjustment ends here


Chris pointed to the \texttt{compute\_blended\_likelihood} function. "This is where the magic happens. \(\alpha\) adjusts the weighting between empirical and theoretical models, and the normalization ensures the probabilities sum to one."

Emmi nodded, following along. "So, the \(\lambda\) parameter is estimated from the average gap size in the training data?"

"Exactly," Chris replied. "And this function dynamically combines both models to create a predictive posterior distribution. The normalization step is crucial—it ensures the probabilities are meaningful and comparable across different scenarios."

She duly showed intrigue nodding ever so perceptably. “So,  your training window is 5000 primes, and you’re predicting the next 100 gaps. Would the model’s predictions evolve during those 100, or is it stuck with what it learned from the training set?”

“It evolves,” Chris said, “but the pace depends on your blending factor and how much the new data deviates from the training data. If gaps follow the expected distribution, the posterior updates are minor. But if something unusual happens, like say a cluster of unexpected large price drops it shifts more dramatically.”

Emmi nodded, her finance background kicking in. “In my world, we’d say this is a classic case of over-fitting. Too many dials to turn, too much room to explain the past without a robust prediction of the future. A good risk model like the ones we try to use at BlackRock, has as few free parameters as possible. Simplicity, of course but naturally limtied as the execs are they are a little more pompous about it in saying that parsimony often wins.”

Chris smiled. “Exactly. That’s the holy grail of modeling: finding the least number of free parameters that still capture the essence of the system. With primes, as with price tick data the gaps don't occur deterministic—they’re inherently random in some ways, but there’s some structure buried in the randomness.”

Emmi, "yes I guess we could argue at a stretch that there are causal reasons for their particular randomness."

Chris nodded-thinking best not to go meta, instead opted to stay grounded and gestured to the plot. “Take this, for instance. The smoothing factor at 0.9 gives a nice balance between short-term volatility and long-term trends. But if I dialed it down to, say, 0.5, the model would overreact to every little blip in the data. And the blending factor \(\alpha\)—too high, and the model relies too much on the empirical gaps; too low, and it’s just the theoretical Poisson model dressed up.”

Emmi frowned. “So what’s the takeaway here? Is this even useful, or is it just mathematical play?”

Chris laughed. “It’s a bit of both. The Bayesian framework is useful because it gives us a probabilistic way to blend empirical data and theory. But if we’re serious about understanding the gaps, we need a model that isn’t so parameter-heavy. Something with deeper constraints, maybe even tied directly to the primes’ intrinsic properties.”

Emmi tilted her head, thinking but not thinking, start a sentence and see where it runs she told herself. “What if we could use a moving posterior window to iteratively refine the model? Like adjusting the priors based on recent data but keeping the framework simple—maybe a single predictive parameter tied to the average gap behavior?”

Chris’s eyes lit up. “Now that’s an idea. Instead of juggling three or four dials, let’s anchor it to something fundamental, like the logarithmic density of primes or a single probabilistic decay factor. It’s less about perfect predictions and more about capturing the underlying trend.”

“And,” Emmi added, with Emily’s advice echoing in her mind, “those patterns can be framed in terms of probabilities. The discovery of irregular primes could be treated as an event with a prior probability. A predictive model can adjust its hit rate dynamically to reflect how well it’s capturing the gaps.”

Chris nodded, his respect for Emmi growing. “Exactly. It’s not just about fitting; rather it’s about understanding the stochastic nature of the process and validating the model iteratively. We’re not chasing numbers; we’re modeling a living, breathing number line.”

He leaned closer, his tone conspiratorial. “But here’s the rub, Emmi. Even with all this, we’re still blending hindsight into foresight. It’s always going to be part art, part science. That’s the nature of primes given they defy pure regularity.”

Emmi smiled. “And maybe that’s why they’re so fascinating. There’s always another layer, another challenge. But if we can even glimpse a little clarity through the noise, maybe that’s enough.”

Chris nodded, raising his now cold cup of tea; again with instant regret. “To models with fewer dials and more insight.” Emmi clinked hers clumsily against his. her eyes perceptibly if fleetingly, locking with his.

% Chris grinned. “You’re thinking like a risk manager now. Validation is all about the hit rate. You compare predicted gaps to actual gaps and measure the fraction of correct predictions. But here’s the twist: accuracy isn’t everything. You also want consistency and robustness. That’s where moving windows help.”

% “Moving windows?” Emmi raised an eyebrow.

% “Picture this,” Chris said, gesturing with his fork. “You divide the data into overlapping segments. For each window, train the model, make predictions for the next segment, and compare. This rolling validation simulates real-world adaptability.”


% Emmi smiled. “You’re still too artful, Chris. But I’ll give you this—it’s starting to make sense.”

\section*{Sisters have a Bayesian Exchange}

\lettrine{E}{mily} stared at the notebook in front of her, its pages filled with calculations on irregular primes. Her fascination with primes — especially those that didn’t divide the numerators of Bernoulli numbers—was a lifelong obsession. Emmi, her younger sister, had brought her coffee, hoping to steal a moment to cement her own understanding of Bayesian ideas. Though their fields couldn’t seem further apart—number theory for Emily, and finance for Emmi—there was a subtle thread that tied their worlds together. Today, Emmi hoped to trade some insights from risk management for clarity on the Bayesian framework.

"Alright, Emmi," Emily said, glancing up. "What’s got you so eager today?"

"Well," Emmi began cautiously, "you’ve been working on irregular primes, right? I’ve been thinking about how we might treat their discovery as an event with prior probability—kind of like breaches in a risk model. I wanted to see if Bayesian methods could fit your work, and, selfishly, I need your help to make sense of them."

\subsection*{From Prime Gaps to Risk Models}


Emily smiled. "Go on. You’ve got my attention."
Emmi leant in. "In finance, we don’t just look at prices—we focus on differences, like percentage or log differences. It’s the changes, not the levels, that matter. Isn’t that a bit like how you study gaps between primes, the first differences? Instead of focusing on individual primes, you’re looking at how they space out as numbers grow larger."

Emily nodded. "Exactly. The gaps are where the real mysteries lie—why do they grow irregularly? And why do some primes, like the irregular ones, seem to resist neat patterns? Keep going."

"In risk management," Emmi continued, "we predict breaches—times when a portfolio’s value falls below a threshold. We assign a prior probability to these events based on historical data, and then we update that probability as new data comes in. For example, if the gaps between primes are like breaches, you could use a Bayesian model to refine your predictions as you test more numbers."

Emily leaned back, intrigued. "So you’re saying the gaps between primes, particularly for irregular primes, could be treated probabilistically? Like they have a prior and posterior probability?"

"Exactly," Emmi replied. "Your prior reflects what you assume about prime gaps—maybe something related to the density of primes or known patterns. Each new gap you calculate refines your posterior, bringing you closer to the truth."

\subsection*{The Square Root Law and Scaling}

% Emily’s eyes lit up. "And where does the square root law fit into this?"

"In finance," Emmi pressing on unusually emboldened in her elder seister's presence, "we use the square root of time rule to model deviations in time series. It’s like a risk model’s built-in scaling law: the deviations grow with the square root of time. Couldn’t the same apply to prime gaps? As \(N\) grows, the variance of the cumulative sum of the gaps—like the sum of the Möbius function—might also grow with \(\sqrt{N}\). It’s a statistical law that seems to persist across different contexts."

Emily tapped her pen against the table. "You might be onto something. If we think of the gaps between primes as events in a sequence, their scaling behavior could follow this law. And it would mean the irregular primes, like outliers in a portfolio, are rare but statistically predictable."

\subsection*{Bayesian Framing for Primes}

Emmi saw her opening. "So here’s what I’m thinking: you assign a prior probability \(P(\text{irregular prime})\) based on your knowledge of irregular primes—like their distribution among smaller numbers. Then, every time you test a number, you update that probability using Bayes’ theorem:
\begin{align*}  
P(\text{prime} \mid \text{evidence}) 
&= \frac{P(\text{evidence} \mid \text{prime}) P(\text{prime})}
      {P(\text{evidence})}.
\end{align*}

The evidence is whether or not the number divides the numerator of a Bernoulli number. Over time, your posterior probability tells you how likely the next number is to be irregular."

Emily grinned. "It’s elegant. And it ties beautifully to what I’ve been working on. Irregular primes are rare, but their discovery could follow a predictive pattern. Bayesian methods would let me refine my search as I gather more data."

"And I get something out of this too," Emmi added quickly. "I get a deeper understanding of Bayesian principles—and the chance to see how the square root law applies beyond finance. I need this, Emily. If I can predict breaches, why not irregular primes?"

Emily chuckled. "Alright, Emmi. You’ve got me thinking differently about primes, and I’ll admit, you’re starting to sound like an academic. But don’t let Chris hear you talking about Bayesian methods with me. He’ll think I’ve stolen his student."

Emmi blushed, her thoughts flickering to Chris. "Well, let’s just say I have a good teacher—and not just you."

As they left the café, Emily felt inspired to revisit her calculations with a Bayesian perspective, while Emmi walked away with a growing confidence in her understanding of risk and probability. Together, they were bridging their worlds, one conversation at a time.

% >>>>>> ANNEX C (cut-sheet v2): the Bayesian blended-likelihood postscript — block commented out below, pointer left in text <<<<<<
\textit{(The Bayesian blended-likelihood postscript --- the full working is set out in Annexe C.)}

% \section*{Postscript}
% The \textbf{Sieve of Eratosthenes} is used to generate prime numbers up to a specified limit:
% \[
% \text{primes} = \{p \mid p \text{ is prime and } p < n\}.
% \]
% The primes are split into:
% \begin{itemize}
%     \item \textbf{Training set:} First $w$ primes for training.
%     \item \textbf{Test set:} Next $p$ primes for testing.
% \end{itemize}
% The gaps between consecutive primes are calculated as:
% \[
% \text{Gap}(i) = \text{Prime}(i+1) - \text{Prime}(i).
% \]
% This yields:
% \begin{itemize}
%     \item $\text{train\_gaps}$ from the training primes.
%     \item $\text{test\_gaps}$ from the test primes.
% \end{itemize}
% The empirical likelihood is computed from the training gaps:
% \[
% L_{\text{empirical}}(g) = \frac{\text{Count of gap } g \text{ in } \text{train\_gaps}}{\text{Total gaps in train\_gaps}}.
% \]
% The theoretical likelihood is based on a \textbf{Poisson distribution}, which models the expected frequency of gaps:
% \[
% L_{\text{theoretical}}(g) = \frac{\lambda^g e^{-\lambda}}{g!},
% \]
% where:
% \begin{itemize}
%     \item $\lambda = \log(\bar{x})$, the mean gap size estimated from the average prime in the training set.
% \end{itemize}
% The two likelihoods are blended as:
% \[
% L_{\text{blended}}(g) = \alpha L_{\text{empirical}}(g) + (1 - \alpha) L_{\text{theoretical}}(g),
% \]
% where $\alpha$ is a blending factor that adjusts the weight of empirical vs. theoretical data. The blended likelihood is normalized:
% \[
% L_{\text{blended}}(g) \gets \frac{L_{\text{blended}}(g)}{\sum_{g} L_{\text{blended}}(g)}.
% \]
% The posterior is updated using the blended likelihood and the prior:
% \[
% \text{Posterior}(g) \propto \text{Prior}(g) \cdot L_{\text{blended}}(g).
% \]
% The prior is initially uniform:
% \[
% \text{Prior}(g) = \frac{1}{|\text{Unique Gaps}|}.
% \]
% Predicted gaps are sampled probabilistically from the posterior distribution:
% \[
% \text{Predicted Gap} \sim \text{Posterior}(g).
% \]
% The posterior is updated incrementally with each observed gap in the test set. The \textbf{hit rate} evaluates the prediction accuracy:
% \[
% \text{Hit Rate} = \frac{\text{Number of Correct Predictions}}{\text{Total Test Gaps}}.
% \]
% 
% This blended approach leverages:
% \begin{itemize}
%     \item \textbf{Empirical Data:} For local accuracy based on observed gaps.
%     \item \textbf{Theoretical Knowledge:} For generalization based on the Poisson distribution.
% \end{itemize}
% 
% By combining these, the model achieves robust predictions for prime gaps across both the training and test sets.
% 
% \newpage
% >>>>>> end of annexed block <<<<<<
\section*{Chris was Stochastic Fair game}

We are in a  quiet pub, Chris and Willard are playing a game. A low table, a backgammon board unfolded,  
coffee cooling beside it.  
Willard is arranging the pieces with calm precision;  
Chris observes quietly nervous, being quietly clueless of what is to be fully unfolding.

\medskip
\textbf{Chris:}  
I’ve never though of playing this sort of thing.  
Why backgammon, Willard?  
There is too much dice involvement to make this skillful surely?

\textbf{Willard:}  
No, its probabilistic but with teeth.  
The board is a field of \emph{affordances} in which
each configuration invites or forbids a class of actions. A good position makes many dice combinations usable;  
a bad one, only a few.  
You could say the position defines the \emph{copula}
between chance and choice. Backgammon is a stochastic, perfect-information game: a deterministic decision process seeded by chance, whose \emph{equity} surface is fully defined in expectation, though never in outcome.

\textbf{Chris:}  
Ok, but I am not sure that helps. What do I have to do to win?

\textbf{Willard:}  
Each player moves their fifteen checkers around the board in opposite directions: White anti-clockwise, Black clockwise, according to the roll of two dice. Checkers may move to open points, to those occupied by their own side, or to single opposing checkers, which are hit and placed on the central bar to re-enter from the start. The goal is to bring all checkers into one’s home board and \emph{bear} them off; the first to remove all wins.

\begin{figure}[h]
    \centering
    \includegraphics[width=0.70\textwidth]{Illustrations/backGammonOpening.png}
    \caption{The opening \href{https://drive.google.com/file/d/1296Id_t4w9Q9d3WY-1NdIF-UFbgU2OOu/view?usp=sharing}{checkers} on a backgammon board.}
    \label{fig:usBoardGammon}
\end{figure}

\textbf{Willard:}  
Because isolated checkers (“blots”) are vulnerable, some dice combinations are far stronger than others at different stages of play. On the opening roll, the best combinations—such as 3-1, 4-2, or 6-1 make key inner-board points and build structure. The weakest openings, like 5-4 or 6-5, move pieces but create exposure without securing points.

\begin{tcolorbox}[colback=blue!5!white, colframe=blue!50!black, title=Backgammon as a Fair Game]
Backgammon can proceed through distinct phases, each with its own relationship between \emph{luck} (dice) and \emph{skill} (position/state):

\scalebox{0.7}{
\begin{center}
\begin{tabular}{llll}
\toprule
\textbf{Phase} & \textbf{Copula Type} & \textbf{Tail Behaviour} & \textbf{Strategic Meaning} \\
\midrule
Opening & Mixed & Balanced & Structure-building; preparing for luck. \\
Blitz / Attack & Gumbel & Upper-tail & Big rolls decide; aggression rewarded. \\
Back-game & Clayton & Lower-tail & Small rolls preserve; timing matters. \\
Bear-off & Gaussian & Symmetric & Pure execution; speed dominates. \\
\bottomrule
\end{tabular}
\end{center}
}
\end{tcolorbox}

\begin{figure}[h]
\centering
\includegraphics[width=0.65\textwidth]{Illustrations/copula_backgammon_phases.png}
\caption{\href{https://colab.research.google.com/drive/1y5HcRJcInjEOKkUlLqoFyBHshgSNH5MT}{Simulated} copula patterns for four backgammon phases. 
Each point represents one hypothetical roll (U) and its resulting positional value (V).}
\end{figure}

\noindent
Each panel shows simulated dependence between two  variables:
$U$ (``roll affordance'' how useful a roll is, structurally) and
$V$ (``equity'' -- how much that roll changes the player's winning chances). 

\paragraph{(1) Opening phase — Mixed copula.}
At the start, players build structure and spread their checkers.
Some rolls are naturally strong (e.g.\ 3--1 makes a key point), others neutral.
Dependence is moderate and uneven: skill arranges the board so that luck can be converted,
but not yet dominated.
This is the \emph{exploratory regime}, where optionality is created and structure remains flexible.

\paragraph{(2) Blitz / Attack phase — Gumbel copula (upper-tail).}
In attacking positions, one player tries to ``close'' the board on an opposing checker.
Good rolls (those that hit and cover) yield large, concentrated gains,
while neutral rolls do little.
Equity clusters in the upper tail: once momentum is seized, it locks in.
This is the \emph{breakout regime}, where a single shock defines the outcome.

\paragraph{(3) Holding / Back-game phase — Clayton copula (lower-tail).}
Both sides anchor defensively, waiting for a shot.
Small, flexible rolls are safe; large, forced rolls can ruin timing.
Dependence strengthens in the lower tail—bad rolls fail together.
This is the \emph{stress regime}, marked by downside correlation and fragile timing.

\paragraph{(4) Bear-off phase — Gaussian copula (linear).}
Once all checkers are home, outcomes depend almost linearly on the dice:
bigger numbers simply remove checkers faster.
Equity and roll quality move in near-symmetric lockstep.
This is the \emph{execution regime}, where play becomes mechanical and variance benign.

\vspace{0.2em}
\textbf{Chris:}  
So the state of the board is a framing of the joint distribution of possible futures\footnote{Gaussian when outcomes mirror luck, Gumbel when good fortune snowballs,
Clayton when bad fortune compounds, and Mixed when both coexist.}  ?

\textbf{Willard:}  
Exactly.   Every roll samples that distribution.   
A point built or blotted changes the geometry of dependence itself.

\textbf{Chris:}  
Then some positions must have more “tail exposure” than others.

\textbf{Willard:}  
Yes. In a {\em blitz}, fortune amplifies—upper-tail Gumbel.  In a \emph{holding game} regime, misfortune compounds—lower-tail Clayton.  A balanced opening? Mixed copula.  And the endgame, when skill no longer matters- bring me doubles, the regime is  pure Gaussian: linear, clean, mechanical.

\textbf{Chris:}  
Ok you are making it sound like volatility regimes in miniature.  Between the dice and the one opponent randomness shape-shifts.

\textbf{Willard:}  
Precisely.  
The board's layout \emph{affords} certain kinds of luck.

The game is now in progress.  Half the board is crowded; two blots exposed.  Chris rolls a four–two.

\vspace{0.2em}
\textbf{Willard:}  
All right. You could move the four from here—(he taps the 13-point)—
down to your 9-point, or bring the two in from the 24-point.
What does each move afford you?

\textbf{Chris:}  
Both are legal, but the first looks safer.
The second leaves a blot.

\textbf{Willard:}  
Exactly, but look deeper.
If you play safe, most of your next rolls will still do nothing.
You’re closing doors ... reducing the number of dice that help.
If you risk the blot, you open space:
more future rolls connect; the board begins to breathe.

\textbf{Chris:}  
So the blot increases my variance.

\textbf{Willard:}  
Variance is only a shadow.
What you’re changing is the \emph{affordance surface}:
the mapping between dice and action.
Before, only a few rolls had value.
Now, many rolls afford something useful.
You’ve reshaped the dependence between chance and outcome.

\textbf{Chris:}  
Like changing the copula?

\textbf{Willard:}  
Yes.  
Every checker moved alters the joint density of luck and equity.
It’s a physical form of modelling.
Each position defines a different joint law.

\textbf{Chris:}  
In markets, we’d call that re-hedging.

\textbf{Willard:}  
Indeed.
A trader re-hedges to change sensitivity;  
a player re-positions to change affordance.
Both are ways of sculpting curvature.

\textbf{Chris:}  
Curvature of what?

\textbf{Willard:}  
Of response.
Convexity if you prefer your own dialect:
the measure of how kindly or cruelly the world reacts to small shocks.
When your board—or your portfolio—curves in your favour,
ordinary randomness begins to work for you.

\textbf{Chris:}  
And if it curves against me?

\textbf{Willard:}  
Then you are playing the wrong game,
and you will be Claytoned before you even realise it.


So after a long backgammon match Chris, has just lost miserably to Willard, who proceeds unhelpfully to use the board to explain what happened to Chris trading model.

\noindent\textbf{Willard:} ``You weren’t beaten by bad rolls, Chris. You were \textit{out of phase}.''

\noindent\textbf{Chris:} ``I was running a lattice-based model. Self-correcting. It tracked curvature in implied volatility. It shouldn’t have broken.''

\noindent\textbf{Willard:} ``Ah, but the market is a whole of players. Once you started dictating prices, the lattice became visible—and the others began playing \textit{you}.''

\noindent\textbf{Chris:} ``They attacked my convexity.''

\noindent\textbf{Willard:} ``Exactly. You built a topology, they found the gradient. In backgammon, that’s the \textit{blitz}. The moment one side smells your structure, they collapse it before it can mature.''

\vspace{0.5em}

\noindent\textbf{Chris:} ``So regime change isn’t something you model, rather it’s something that \textit{happens to you}.''

\noindent\textbf{Willard:} ``Yes. Every model lives only within its phase of validity. Then comes a roll that doesn’t fit the prior. That’s the moment of Gumbel revelation, when an upper-tail shock hits, when markets mass psychology realigns.''

\noindent\textbf{Chris:} ``And so liquidity evaporates, and even the simplest bi-variate correlations, become ...''

\noindent\textbf{Willard:} ``Claytonian. Lower-tail lockstep as everyone is running for the same door. You’d built a perfect lattice under a shifting sky. Convexity was supposed to protect you, but it marked you instead. They could see your structure such that every rescue trade was a signal flare.''

\vspace{0.5em}

Willard now annoyingly tapping the board, on a roll with little resistance.  
``These phases changes, the regime changes they came too fast. You started as Gaussian—balanced, diffuse, ...  and then your success made you Gumbel, all upper-tail ambition.  That turned you Clayton, well it should have done, should have forced you to be defensive. And when you belatedly tried to bear off, to exit gracefully, there was no linearity left—no Gaussian calm to return to.  The market had changed state around you.''

\vspace{0.5em}

\noindent\textbf{Chris:} ``So as in this silly game equilibrium isn’t safety ... it’s just a phase between aggressions.''

\noindent\textbf{Willard:} ``i am sorry that this is true. Every strategy works—until it redefines the surface it’s standing on.''

\vspace{1em}

Chris looks forlornly down at his last two checkers marooned on the bar of the last game his. Willard has already cleared his board. The cube stands at four, and the loss given it's a backgammon is $3\times4=$ Twelve points swing across the table. Another ritual humiliation.

\vspace{1em}

\subsection*{Epilogue: Affordance and Curvature}

\textbf{Chris:}    
In bond maths we say a portfolio is convex when small shifts in rates produce asymmetric gains: you make more when yields fall than you lose when they rise.  Backgammon feels like that:  some positions have positive convexity, others negative.

\textbf{Willard:}  
Yes that works as an analogy.  Convexity in rates is curvature in price–yield space.  In backgammon I guess it’s curvature in the equity–luck surface.

\textbf{Chris:}  
So when you blitzed me earlier ... you’d built positive convexity:  a handful of high rolls gave you explosive returns.

\textbf{Willard:} 
Yes Each copula family corresponds to a different way that \emph{random events}
(luck) and \emph{system state} (skill or structure) combine to produce outcomes; convexity describes how the outcome surface bends under perturbation. That was Gumbel copula territory.  Then, when you played the back game, you built negative convexity: bad rolls punished you more than good ones helped. Clayton.

\medskip
\textbf{Chris:}  
I used to think the market afforded opportunity. That is,  that convexity was a door I could open wider just by being clever.  Now I see it was a door that closes faster  the more people reach for the handle.

\textbf{Willard:}  
That’s an old story.   In games the geometry is fixed;  in markets it moves beneath you.   The more you exploit an affordance, the more it disappears.  Liquidity is a receding surface.

\textbf{Chris:}  
So I built my book like a backgammon position in trying to make every roll work for me.   But when the big funds turned,  
they changed the copula itself.   My convexity flattened overnight.

\textbf{Willard:}  
You were Claytoned, not Gumbled.  
Lower-tail dependence, in which every bad move was reinforcing the next.  It happens when the structure stops affording you any escape.

\textbf{Chris:}  
And there’s no Gaussian phase in markets?

\textbf{Willard:}  
Only in hindsight.  
Markets never bear off.  
They just reset the board and ask again what your structure affords.

\textbf{Chris:}  
So in the end, the difference between a game and a market   is that one ends in closure,  and the other keeps reconfiguring its own geometry. The memory of that wooden board will linger

\textbf{Willard:}  
Yes.  And the skill ... whether playing dice or derivatives ... is knowing when your convexity is real,  and when it’s just a shape the market allows you to see for a moment.

\noindent\footnotesize{\textcolor{gray}{%
\textbf{Note.} ``Regime change'' as in a phase change is defined by its dependence geometry: the \em{copula} linking variables.  When external feedback loops (e.g., reflexivity in Soros’ sense) alter those dependencies, the system crosses a phase boundary.  Markets and games alike become \textit{endogenous}: the act of modeling changes what is modeled.}}

% >>>>>> ANNEX C (cut-sheet v2): copulas and the Betti dynamics worked example — block commented out below, pointer left in text <<<<<<
\textit{(Copulas and the Betti dynamics worked example --- the full working is set out in Annexe C.)}

% \subsection*{Copulas: The Geometry Beneath Correlation}
% 
% A \textbf{copula} isolates the structure of dependence between random variables
% from their individual behaviour.  
% For any joint distribution \(F_{XY}(x,y)\) with continuous marginals \(F_X, F_Y\),
% there exists a copula \(C:[0,1]^2 \to [0,1]\) such that
% \[
% F_{XY}(x,y) = C(F_X(x),\,F_Y(y)).
% \]
% By defining \(U=F_X(X)\) and \(V=F_Y(Y)\),
% we obtain two uniform variables on \([0,1]\) whose joint law is \(C(u,v)\).
% 
% \emph{Sklar’s theorem} tells us:  
% copulas are the invariant core of dependence, independent of marginal shape.
% 
% If \((X,Y)\) are jointly normal with correlation \(\rho\),
% their copula is the \textbf{Gaussian copula}:
% \[
% C_\rho(u,v) = 
% \Phi_\rho\!\bigl(\Phi^{-1}(u),\,\Phi^{-1}(v)\bigr),
% \]
% where \(\Phi^{-1}\) is the univariate normal quantile function
% and \(\Phi_\rho\) is the bivariate normal CDF with correlation \(\rho\).
% 
% \begin{quote}
% Correlation measures how two variables line up on average.\\
% Copulas reveal \emph{where} they align---in the middle, or in the tails.
% \end{quote}
% 
% The familiar Pearson coefficient \(\rho\) is thus simply the coordinate of the Gaussian copula. thus not the dependence itself, rather merely the parameter of this one copula family describing the \emph{elliptical} symmetry of Linear correlation. Other copulas reveal shapes of dependence that correlation alone cannot express:
% 
% \begin{itemize}
%   \item \textbf{Clayton copula} --- strong \emph{lower-tail} connection:
%         bad outcomes cluster together.
%   \item \textbf{Gumbel copula} --- strong \emph{upper-tail} connection:
%         extreme good outcomes reinforce each other.
%   \item \textbf{Frank copula} --- balanced but non-elliptical symmetry.
% \end{itemize}
% 
% So dependence can have curvature, asymmetry, or saturation
% independent of how their marginal distributions scale. In games, markets, or physical systems, the copula tells us whether luck and skill reinforce each other in the tails, cancel symmetrically, or diverge. Correlation is only the linear shadow of that richer geometry.
% 
% \begin{center}
% \begin{tikzpicture}[scale=1.1]
%   % Axes helper
%   \def\axes{
%     \draw[->,thin,gray] (0,0)--(1.1,0) node[right,gray]{$U$};
%     \draw[->,thin,gray] (0,0)--(0,1.1) node[above,gray]{$V$};
%     \draw[gray!40] (0,0) rectangle (1,1);
%   }
%   % Gaussian
%   \begin{scope}[shift={(0,0)}]
%     \axes
%     \draw[blue!60, thick, domain=0:1,samples=30] plot (\x,\x+0.05*rnd);
%     \node[below right,blue!70!black]{\small Gaussian};
%   \end{scope}
%   % Gumbel
%   \begin{scope}[shift={(2.5,0)}]
%     \axes
%     \draw[red!70, thick, domain=0:1,samples=30] plot (\x,{\x^(0.8)+0.1*rnd});
%     \fill[red!20,opacity=0.4] (0.6,0.6) rectangle (1,1);
%     \node[below right,red!70!black]{\small Gumbel};
%   \end{scope}
%   % Clayton
%   \begin{scope}[shift={(5,0)}]
%     \axes
%     \draw[orange!80!black, thick, domain=0:1,samples=30] plot (\x,{1-(1-\x)^(0.8)+0.1*rnd});
%     \fill[orange!20,opacity=0.4] (0,0) rectangle (0.4,0.4);
%     \node[below right,orange!80!black]{\small Clayton};
%   \end{scope}
%   % Mixed
%   \begin{scope}[shift={(7.5,0)}]
%     \axes
%     \draw[violet!70, thick, domain=0:1,samples=30] plot (\x,{0.9*\x+0.1*sin(5*deg(\x))+0.05*rnd});
%     \node[below right,violet!70!black]{\small Mixed};
%   \end{scope}
% \end{tikzpicture}
% \end{center}
% 
% \noindent
% Each copula describes a distinct ``shape of dependence'' between two percentile
% variables $(U,V)$. \textbf{Gaussian}: a smooth, symmetric diagonal—good outcomes mirror good rolls linearly. \textbf{Gumbel}: clustering in the \emph{upper-right}—excellent rolls amplify success.\textbf{Clayton}: clustering in the \emph{lower-left}—poor rolls amplify loss. \textbf{Mixed}: a hybrid with moderate, uneven dependence.
% \newpage
% \subsection*{Example: Betti Dynamics on a 4×4 Mini–Checkers Lattice}
% 
% We illustrate the homological interpretation of play using a reduced 4×4 checkers board 
% containing only the dark squares
% \[
% S = \{(1,2),(1,4),(2,1),(2,3),(3,2),(3,4),(4,1),(4,3)\},
% \]
% with rows numbered from bottom ($1$) to top ($4$).
% White moves diagonally upward.
% 
% \paragraph{Step 0 (initial position).}
% White occupies $(1,2)$ and $(1,4)$, Black occupies $(4,1)$ and $(4,3)$.
% For White, legal forward steps are
% \[
% (1,2)\to\{(2,1),(2,3)\}, \qquad (1,4)\to(2,3),
% \]
% giving the vertex set
% \[
% V_0=\{(1,2),(1,4),(2,1),(2,3)\},
% \quad 
% E_0=\{(1,2)\!-\!(2,1),(1,2)\!-\!(2,3),(1,4)\!-\!(2,3)\}.
% \]
% Since all vertices are connected, the mobility graph $G_0$ has
% \[
% b_0(G_0)=1, \qquad b_1(G_0)=0.
% \]
% White’s mobility forms one cohesive region without loops.
% 
% \paragraph{Step 1 (White moves the middle piece).}
% White advances $(1,2)\to(2,1)$, leaving
% \[
% V_1=\{(2,1),(1,4),(3,2),(2,3)\}, \quad
% E_1=\{(2,1)\!-\!(3,2),(1,4)\!-\!(2,3)\}.
% \]
% The two pairs are disconnected, giving
% \[
% b_0(G_1)=2,\qquad b_1(G_1)=0.
% \]
% The mobility surface fragments: White now plays in two disjoint corridors.
% 
% \paragraph{Step 2 (Black replies).}
% Black advances $(4,3)\to(3,2)$, blocking one lane but enabling a jump:
% \[
% (2,1)\to(4,3) \quad \text{over the black at } (3,2).
% \]
% Adding this long ``jump edge'' extends the reach of $(2,1)$ but keeps components disjoint for now:
% \[
% b_0(G_2)=2,\qquad b_1(G_2)=0.
% \]
% The topology is poised for change—an impending capture corresponds to a collapsing boundary.
% 
% \paragraph{Step 3 (White executes the jump).}
% White captures $(3,2)$ and lands on $(4,3)$, merging the reachable regions:
% \[
% b_0(G_3)=1,\qquad b_1(G_3)=0.
% \]
% A local topological singularity has been resolved; the mobility space becomes connected again.
% 
% \paragraph{Interpretation.}
% The Betti sequence
% \[
% (b_0,b_1):\ (1,0)\ \rightarrow\ (2,0)\ \rightarrow\ (2,0)\ \rightarrow\ (1,0)
% \]
% tracks the fragmentation and re–coalescence of the player’s decision space.
% An increase in $b_0$ marks regime bifurcation (loss of coherence);
% a decrease signals consolidation (capture or merger).
% A future \emph{king} would introduce backward mobility, closing a loop
% and producing the first nonzero $b_1$.
% 
% \paragraph{Willard’s note.}
% \begin{quote}\small
% ``Each checker move reshapes the topology.
% When $b_0$ splits, the field fragments; when $b_0$ contracts, it fuses.
% Only when a king appears does $b_1$ breathe—then the board acquires memory.''
% \end{quote}
% 
% \begin{figure}[h!]
% \centering
% \begin{tikzpicture}[scale=1.15, every node/.style={circle, minimum size=6pt, inner sep=0pt}]
% \tikzset{
%     node/.style={circle, draw=black, fill=gray!15, minimum size=10pt},
%     whitep/.style={circle, draw=black, fill=white, thick, minimum size=10pt},
%     blackp/.style={circle, draw=black, fill=black, minimum size=10pt},
%     step/.style={dotted, thick, gray!70},
%     jump/.style={dashed, thick, red!70},
% }
% 
%%% PANEL (a): Start
% \node at (0,2.8) {\textbf{(a) Start: $(b_0,b_1)=(1,0)$}};
% nodes
% \node[node] (a12) at (0,0) {};
% \node[node] (a14) at (1.2,0) {};
% \node[node] (a21) at (-0.6,0.8) {};
% \node[node] (a23) at (0.6,0.8) {};
% pieces
% \node[whitep] at (0,0) {};   % (1,2)
% \node[whitep] at (1.2,0) {}; % (1,4)
% edges
% \draw[step] (0,0) -- (-0.6,0.8);
% \draw[step] (0,0) -- (0.6,0.8);
% \draw[step] (1.2,0) -- (0.6,0.8);
% 
%%% PANEL (b): After White moves (fragmented)
% \node at (4.0,2.8) {\textbf{(b) After move: $(b_0,b_1)=(2,0)$}};
% \node[node] (b21) at (3.4,0.8) {};
% \node[node] (b32) at (3.4,1.6) {};
% \node[node] (b14) at (5.0,0) {};
% \node[node] (b23) at (4.4,0.8) {};
% pieces
% \node[whitep] at (3.4,0.8) {}; % (2,1)
% \node[whitep] at (5.0,0) {};   % (1,4)
% edges
% \draw[step] (3.4,0.8) -- (3.4,1.6);
% \draw[step] (5.0,0) -- (4.4,0.8);
% 
%%% PANEL (c): Black replies → jump path reconnects
% \node at (8.0,2.8) {\textbf{(c) Jump phase: $(b_0,b_1)=(1,0)$}};
% nodes
% \node[node] (c21) at (7.4,0.8) {};
% \node[node] (c14) at (9.0,0) {};
% \node[node] (c23) at (8.4,0.8) {};
% \node[node] (c32) at (7.4,1.6) {};
% \node[node] (c43) at (8.6,2.4) {};
% pieces
% \node[whitep] at (7.4,0.8) {}; % (2,1)
% \node[whitep] at (9.0,0) {};   % (1,4)
% \node[blackp] at (7.4,1.6) {}; % (3,2) black moved down
% edges
% \draw[step] (9.0,0) -- (8.4,0.8);
% \draw[jump] (7.4,0.8) -- (8.6,2.4); % jump edge over (3,2)
% 
% \end{tikzpicture}
% \caption{Evolution of White’s mobility graph on the 4×4 mini–checkers lattice.  
% (a) Initial position: a single connected mobility component $(b_0,b_1)=(1,0)$.  
% (b) After White moves, the graph fragments into two independent regions $(b_0=2)$.  
% (c) Black advances, enabling a long jump edge (red dashed), which re-connects the graph $(b_0=1)$.  
% At no stage does a loop appear ($b_1=0$), but a future king could generate $b_1>0$.}
% \end{figure}
% 
% 
% 
% \newpage
% 
% >>>>>> end of annexed block <<<<<<
\subsection*{Metrics and Monetization}

\lettrine{E}{leanor’s} first stop was New York City. The boardroom  was sleek, filled with steely-eyed investors more interested in quarterly returns than abstract theorems. Eleanor presented the Merriweather Circle’s groundbreaking work—but her words bounced off the polished walls.

\begin{figure}[h]
    \centering
    \includegraphics[width=\textwidth]{Illustrations/vign-mixedBoardrooms.png}
    \caption{Steely-eyed investors and sharp suits.}
    \label{fig:usBoardroom}
\end{figure}

A sharp suitor leaned forward, his voice dripping with skepticism. \textit{“Professor Merriweather, this is fascinating, but how does it scale? Where’s the application? Gene therapy, quantum computing, something tangible?”}
\bigskip

Eleanor straightened. \textit{“Mathematics is the bedrock of every technological revolution. Theorems proven today become the algorithms of tomorrow. Without the groundwork, there are no towers.”}
\bigskip

While her conviction was clear, the words fell flat in the ROI-obsessed room. A different voice chimed in from the corner, summing up the group’s sentiments. \textit{“We fund what solves problems, not what creates them.”}
\bigskip

As she left, Eleanor couldn’t help but feel the sting of their misunderstanding. These people saw mathematics as a tool, never an end in itself. Still, she resolved to refine her pitch—it wasn’t just about changing their minds; it was about preserving the integrity of her Circle.

\subsection*{Skepticism and Eccentricity}

Her next stop in London felt different but no less frustrating. At a pharmaceutical conglomerate, the executives were polite, sipping their teas with an air of bemusement. Eleanor outlined the Circle’s mission, leaning on Tiziana and Chris Everett’s work in wavelets and Fourier analysis.

\smallskip

\textit{“How can we integrate your ideas into oour big data machine learning program?”} a director asked, his accent clipped and precise. \textit{“Surely there’s some tie-in?”}

Eleanor faltered slightly but recovered. \textit{“Wavelets and Fourier transforms could absolutely enhance your AI models, particularly in identifying patterns in noisy datasets.”}
\bigskip

There were polite nods, but Eleanor could tell it wasn’t enough. As the meeting ended, an older woman approached her. \textit{“Refine your pitch, dear. They won’t understand unless you connect it to something they already do.”}
\smallskip

Eleanor walked out of the building with a growing awareness: the purity of her Circle’s work was its greatest strength, but it might also be its greatest liability.

\subsection*{Larissa Volkmann}

By the time Eleanor reached Berlin, she was weary but determined. The boardroom  was dominated by Larissa Volkmann. A towering figure, Larissa commanded the room..

\begin{figure}[h]
    \centering
    \includegraphics[width=1.0\textwidth]{Illustrations/vign-larissaBoardroom.png}
    \caption{Larissa Volkmann’s boardroom: a maverick’s playground.}
    \label{fig:larissaBoardroom}
\end{figure}

Eleanor began, but Larissa interrupted her with a booming laugh. \textit{“Enough about purity and ideals! Tell me how this could work for me.”}

\smallskip

Eleanor hesitated, then pivoted. \textit{“Wavelets and Fourier transforms. Tiziana and Chris Everett’s methods could optimize your data pipelines. And Nambu’s monopole primes? They might hold implications for encryption.”}
\smallskip

Larissa leaned back, grinning. \textit{“Now you’re talking. But I’ll fund the pure work too—on one condition. I want your Circle to get its hands a little dirty with AI models and real-world data. Pure math feeds applied math, Eleanor. Let’s not pretend they’re separate.”}
\smallskip

For the first time, Eleanor felt a spark of hope. Larissa got it—more than anyone else had so far.


\subsection*{Quantum Challenges in Networking}

Eleanor’s fourth meeting, back in London, brought her to a consultancy firm partnered with a telecom conglomerate. Their mission: future-proofing global networks against the rise of quantum computing and the vulnerabilities it posed to traditional encryption.
\smallskip
Dr. Alan Greaves, a seasoned consultant, greeted her. \textit{“Professor Merriweather, we’ve followed the Circle’s work. Help us bridge the gap between theoretical mathematics and practical solutions.”}
\smallskip

Eleanor began, outlining the looming threat posed by Shor’s algorithm and the need for quantum-resilient encryption. \textit{“The days of RSA and ECC encryption are numbered. The Merriweather Circle is exploring lattice-based cryptography and other number-theoretic constructs that could offer alternatives.”}
\smallskip

The room collectively drew closer, \textit{“Rough primes,” Eleanor continued, “primes that avoid small factors, could inspire algorithms resistant to quantum attacks. And our work on Fibonacci-inspired tilings could enhance quantum communication by organizing qubit states with reduced error rates.”}
\smallskip

A young engineer, Priya Malik, asked, \textit{“Are you suggesting this work could aid in secure quantum communication?”}

\begin{figure}[H]
\includegraphics[width=0.8\linewidth]{Illustrations/vign-eleanorCOMpSoence.png}
\caption{Eleanor having a code}
\end{figure}

Eleanor nodded. \textit{“Absolutely. By merging number theory with quantum algorithms, we can design protocols that thrive in a post-quantum world.”}
\bigskip

For the first time, Eleanor felt her Circle’s work resonate deeply with a practical audience. But as the meeting ended, she couldn’t shake the unease creeping in. To what extent could the Circle’s purity be preserved? Was she compromising their ideals to secure their survival?
\bigskip

As Eleanor returned to her hotel, she mused on the whirlwind tour. In New York, she’d faced the cold calculus of monetization. In London, the skeptical pragmatism of pharma. In Berlin, Larissa’s bold pragmatism had offered hope. And with the telecom consultants, she’d seen a glimpse of their work’s real-world impact.
\bigskip

Yet beneath it all lay a growing tension. How could she balance the Circle’s commitment to pure mathematics with the demands of a world increasingly obsessed with applications? Could she hold true to their ideals while navigating the compromises necessary to keep their work alive?
\bigskip

Eleanor sighed, staring out at the city lights. The stakes had never been higher, but neither had the potential. One thing was clear: the Circle’s path forward would require not just brilliance but courage—and Eleanor was ready to lead.



